\section{Exposure and duration laws}
\label{sec:laws}

\subsection{The necessary side}

For $n$ founders of a single type with $\eta=1-q_i>0$ and $c>0$, risk at most
$\delta\in(0,1)$ necessarily requires, by taking the $n$th root in
\eqref{eq:main},
\begin{equation}\label{eq:necessary}
 B\ \ge\ B_{\mathrm{nec}}(n,\delta):=\frac1c
 \left[\log\frac{\eta}{1-(1-\delta)^{1/n}}\right]_+ ,
\end{equation}
with $[x]_+=\max(x,0)$. For $0<\delta\le1/2$,
\begin{equation}\label{eq:rootbounds}
 \frac\delta n\ \le\ 1-(1-\delta)^{1/n}
 \ \le\ \frac{-\log(1-\delta)}n\ \le\ \frac{2\delta}n .
\end{equation}
The first inequality follows from concavity of $x^{1/n}$ at $x=1$, the second
from $1-e^{-y}\le y$, the third by integrating $(1-x)^{-1}\le2$ on $[0,\delta]$.
Hence $B_{\mathrm{nec}}(n,\delta)\ge c^{-1}[\log(\eta n/(2\delta))]_+$.

\subsection{The sufficient side}

\begin{proposition}[Weighted drift]\label{prop:weight}
Let $v_*$ be an admissible constant action and suppose $w>0$ satisfies
$A_{v_*}w\le-\gamma w$ with $\gamma>0$, where $A_{v_*}$ is
\eqref{eq:meanmatrix}. Then under the constant action $v_*$,
\begin{equation}\label{eq:upper}
 \Prob_z(|Z_t|>0)\ \le\ \frac{\sum_iz_iw_i}{w_{\min}}\,e^{-\gamma t}.
\end{equation}
For $n$ identical founders of type $i$ this gives risk at most $\delta$ once
$t\ge\gamma^{-1}\log(Cn/\delta)$ with $C=w_i/w_{\min}$, at exposure
\begin{equation}\label{eq:sufficient}
 B_{\mathrm{suff}}=\frac{v_*}{\gamma}\,\log\frac{Cn}{\delta},
\end{equation}
provided $v_{\max}\ge v_*$ and the horizon is at least
$\gamma^{-1}\log(Cn/\delta)$. Extinction is absorbing, so the bound persists
after any continuation without immigration.
\end{proposition}

\begin{proof}
$W(z)=\sum_iz_iw_i$ has drift $\sum_iz_i(A_{v_*}w)_i\le-\gamma W(z)$.
Localization with the first-moment domination of Section~\ref{sec:model} and
Gronwall give $\E W(Z_t)\le W(z)e^{-\gamma t}$. On $\{|Z_t|>0\}$ we have
$W(Z_t)\ge w_{\min}$, so Markov's inequality gives \eqref{eq:upper}.
Substituting $t=\gamma^{-1}\log(Cn/\delta)$ and $B=v_*t$ gives the rest.
\end{proof}

Since $A_{v}$ is Metzler, $A_vw\le-\gamma w$ for some $w>0$ implies
$\spb(A_v)\le-\gamma$, and conversely a positive witness exists with
$\gamma$ arbitrarily close to $-\spb(A_v)$ when $A_v$ is irreducible; the
two-sided Collatz--Wielandt bounds
\begin{equation}\label{eq:cw}
 \min_i\frac{(Aw)_i}{w_i}\ \le\ \spb(A)\ \le\ \max_i\frac{(Aw)_i}{w_i}
 \qquad(w>0,\ A\ \text{Metzler})
\end{equation}
turn any positive rational vector into a certified spectral bracket
\citep{berman1994}. We use \eqref{eq:cw} throughout
Section~\ref{sec:memory}; a diagonal similarity by $\diag(w)$ preserves the
spectrum and produces a Metzler matrix whose row sums lie between the two
displayed quantities.

\begin{corollary}[Matching logarithmic order]\label{cor:order}
Fix a source with $\eta>0$, $c>0$ as in \eqref{eq:cert} and an admissible
constant action with $\gamma>0$ as in Proposition~\ref{prop:weight}. Define
$B^*(n,\delta)$ as the infimum of pathwise budget bounds among policies
attaining risk at most $\delta$, allowing any finite deterministic horizon;
either the finite-time or the declared eventual endpoint may be used. With the
amplitude fixed so as to admit that action,
\[
 B^*(n,\delta)=\Theta\bigl(\log(n/\delta)\bigr)
 \qquad\text{as }n/\delta\to\infty,\quad0<\delta\le1/2 .
\]
\end{corollary}

\begin{proof}
Equations \eqref{eq:necessary}--\eqref{eq:rootbounds} give
$B^*\ge c^{-1}[\log(\eta n/(2\delta))]_+$ and \eqref{eq:sufficient} gives
$B^*\le(v_*/\gamma)\log(Cn/\delta)$. All constants are fixed independently of
$n$ and $\delta$. The upper bound supplies a finite horizon for each pair. It
makes no feasibility assertion at a fixed short horizon, and it is vacuous when
no admissible constant action has $\gamma>0$ --- precisely the situation
Theorem~\ref{thm:amp} converts into a positive survival floor.
\end{proof}

The two constants $c^{-1}$ and $v_*/\gamma$ need not agree, and closing the gap
between them is a source-specific question. Section~\ref{sec:combo} exhibits an
actuator in the molecular source for which the ratio of the two is below
$1.45$, uniformly in the size of the inherited state.
